\documentclass{article}
\usepackage[UTF8]{ctex}
\usepackage{amsmath}
\usepackage{amssymb}
\usepackage{graphicx}
\usepackage{geometry}
\geometry{a4paper, margin=2cm}

\newcommand{\problemname}[1]{\begin{center}\Large\bfseries #1\end{center}}
\usepackage{enumitem}

\begin{document}
\problemname{温度转换}

花椰妹最近学习了如何将摄氏温度和华氏温度互相转换。

她发现可以用这两个公式来进行转换：

\begin{itemize}
  \item 从输入的摄氏温度转为华氏温度可以通过 $1.8\times\text{输入的温度}+32$ 得到；

  \item 从输入的华氏温度转为摄氏温度可以通过 $5\times(\text{输入的温度}-32)/9$ 得到。

\end{itemize}
她现在需要你帮忙写一个程序，可以快速地完成温度转换计算。

\subsection*{输入格式}
输入的第一行是一个字母，是输入的温度的温标（单位）：可能是 \texttt{C}（输入的温度是摄氏温度）或 \texttt{F}（输入的温度是华氏温度）。

输入的第二行是一个数字，表示要转换的温度值。这个温度值的范围在 $-459.67$ 到 $10000$ 之间。

\subsection*{输出格式}
输出一行，包含一个数字。如果是输入的是摄氏温度，转换后保留小数点后 $2$ 位；如果输入的是华氏温度，转换后保留小数点后 $5$ 位。

\subsection*{样例输入}
\begin{verbatim}
F
41
\end{verbatim}

\subsection*{样例输出}
\begin{verbatim}
5.00000
\end{verbatim}


\end{document}
